\documentclass[12pt, letterpaper]{article}
\usepackage[margin=1.2in]{geometry}
\usepackage{ebgaramond}
\usepackage[T1]{fontenc}
\usepackage{microtype}
\usepackage{setspace}
\usepackage{parskip}
\usepackage{titling}
\usepackage{fancyhdr}

\setlength{\parindent}{2em}
\onehalfspacing

\pagestyle{fancy}
\fancyhf{}
\fancyhead[L]{\small\textit{The Weather Eye}}
\fancyhead[R]{\small\textit{NASA Mission Series v1.15}}
\fancyfoot[C]{\small\thepage}
\renewcommand{\headrulewidth}{0.4pt}

\title{\Huge\textbf{The Weather Eye}\\[0.5em]\large\textit{A short story from Mission 1.15 --- TIROS-1}}
\author{NASA Mission Series\\[0.3em]\small Miles Tiberius Foxglove --- Mukilteo, WA}
\date{March 30, 2026}

\begin{document}
\maketitle
\thispagestyle{empty}

\vspace{1em}
\noindent\rule{\textwidth}{0.4pt}
\vspace{1em}

\noindent The photograph arrived at 0947 Eastern Standard Time on a Tuesday morning in April, and it changed everything I thought I knew about the weather.

I had been a forecaster at the U.S. Weather Bureau for eleven years. I had drawn thousands of surface analysis charts by hand, plotted upper-air soundings on skew-T diagrams, argued with my colleagues about the position of cold fronts based on conflicting station reports from cities hundreds of miles apart. I had predicted rain that never came and missed ice storms that shut down entire states. I was good at my job. I was also, in a fundamental way that I could not have articulated before that Tuesday, blind.

The photograph was grainy. It was black and white, the resolution poor by any photographic standard, the contrast uneven. It had been taken by a television camera --- not a photographic camera, a \emph{television} camera, a vidicon tube designed for broadcasting, repurposed and pointed downward --- from an altitude of seven hundred kilometers, transmitted as a radio signal to a ground station at Fort Monmouth, New Jersey, and printed on a facsimile machine. What it showed was clouds.

Clouds. The word is inadequate. What it showed was a weather system.

\begin{center}
$\bullet\quad\bullet\quad\bullet$
\end{center}

I had never seen a weather system before. I had inferred weather systems. I had deduced their presence from barometric pressure readings and wind directions and temperature gradients and the reports of ship captains and pilots and weather station observers scattered across a continent. I had drawn isobars on charts --- those graceful curving lines connecting points of equal pressure --- and from the pattern of isobars I had imagined the shape of the storm. The tightly packed isobars meant strong winds. The closed isobars meant a low-pressure center. The warm front extended southeast, the cold front trailed southwest, and where they met there was rain, and where the cold front overtook the warm front there was an occluded front, and all of this I knew because thirty-seven weather stations had reported their observations at the same hour, and I had plotted those observations on a chart, and I had connected the dots.

But I had never \emph{seen} the storm. I had never seen the spiral. I had never seen the arms of cloud reaching outward from the center like the arms of a galaxy, curving with the Coriolis force, white against the dark ocean, unmistakable and enormous and real. And here it was, in a grainy photograph taken from space by a drum-shaped satellite the size of a washing machine, spinning at twelve revolutions per minute in its orbit around the Earth, its camera seeing the ground for only a third of each rotation, building up an image one scan line at a time like a television set painting a picture --- and the picture it painted was more information about the weather than I had gathered in a decade of surface observations.

I sat at my desk and stared at the photograph for a long time.

\begin{center}
$\bullet\quad\bullet\quad\bullet$
\end{center}

They called the satellite TIROS --- Television InfraRed Observation Satellite. It had been launched the day before, April 1, 1960, from Cape Canaveral, on a Thor-Able rocket. It weighed a hundred and twenty kilograms. Its body was a drum, forty-two inches in diameter and nineteen inches tall, covered in solar cells that powered the two television cameras and the radio transmitters. It spun on its axis for stability, the way a bullet spins, and the spin meant the cameras could only photograph the Earth during the portion of each revolution when they happened to point downward. Roughly a hundred degrees of arc out of three hundred and sixty. The rest of the time, they photographed the blackness of space.

The cameras were vidicon tubes --- the same technology that made live television possible. A photoconductive surface, antimony trisulfide, changed its electrical resistance in response to the light pattern focused on it by a lens. An electron beam scanned the surface in a raster pattern, reading out the stored image as a varying electrical signal. Five hundred lines of resolution. The wide-angle camera had a field of view of a hundred and four degrees, covering an area roughly fifteen hundred kilometers across from the satellite's altitude. The narrow-angle camera saw a much smaller area --- about a hundred and fifty kilometers --- but in finer detail.

In its first day, TIROS-1 returned dozens of photographs. Within a week, hundreds. By the time its cameras failed after seventy-eight days of operation, it had transmitted twenty-two thousand, nine hundred and fifty-two images of the Earth's cloud cover. It was the most weather data any single instrument had ever collected. And the photographs were not merely data --- they were a new way of seeing.

\begin{center}
$\bullet\quad\bullet\quad\bullet$
\end{center}

The photograph I was staring at showed a typhoon. We had known, from ship reports and island weather stations, that there was a tropical disturbance in the western Pacific. We had estimated its position within a few hundred nautical miles. We had warned the shipping lanes. But we did not know its shape. We did not know its size. We did not know how tightly wound it was, how many arms it had, how far the cirrus shield extended outward from the eye. We could not see the eye.

Now I could see the eye. It was there in the photograph, a dark circle at the center of the spiral, the calm center where descending air cleared the clouds and the sea surface reflected sunlight upward to the camera. The spiral arms extended outward for hundreds of kilometers, curving counterclockwise --- the Northern Hemisphere signature of a cyclonic vortex --- and the cloud bands were layered, with dense convective towers in the inner wall and thinner cirrus streaming outward at the periphery. The entire system was perhaps a thousand kilometers in diameter. It was, from seven hundred kilometers above, stunningly beautiful.

I thought of jellyfish. I do not know why the association came; perhaps it was the pulsing quality of the spiral, the sense of a living thing contracting and expanding, drawing energy from its environment and converting it into organized motion. A moon jellyfish --- \emph{Aurelia aurita} --- is translucent, drum-shaped, and moves by rhythmic contraction of its bell. TIROS-1 was drum-shaped and moved by the rhythmic contraction of orbital mechanics, the constant falling that produces a circle. The jellyfish pulses through the ocean, carried by currents it does not control, sensing the light through simple photoreceptors arrayed around its bell margin. TIROS-1 pulsed through its orbit, carried by gravity, sensing the light through a photoconductive surface behind a lens. Both were simple organisms --- the jellyfish biologically simple, the satellite technologically simple --- that revealed the vast complexity of the medium through which they moved.

\begin{center}
$\bullet\quad\bullet\quad\bullet$
\end{center}

The storm in the photograph was the beginning of a thought that took me years to fully form. The thought was this: we had been forecasting the weather from inside the weather. We had been fish trying to understand the ocean by tasting the water immediately around us. We measured temperature and pressure and humidity and wind at a single point, and from that single point we extrapolated outward, building a mental model of the atmosphere that was always incomplete, always approximate, always limited by the spacing of our observation stations and the speed at which data could be collected and transmitted and plotted. We were blind men describing an elephant by touch --- here a leg, here a trunk, here an ear --- and TIROS-1 had given us eyes.

The eyes were crude. Five hundred lines of resolution, two frames per second, only photographing when the spinning drum pointed the camera earthward. The coverage was limited to the latitude band between forty-eight degrees north and south --- the orbital inclination determined what the satellite could see. The images were distorted at the edges by the wide-angle lens and by the oblique viewing geometry at large look angles. The vidicon tube had a persistence problem: bright clouds left ghost images that faded slowly, contaminating subsequent frames. The tape recorder that stored images for later playback had limited capacity. The satellite itself lasted only seventy-eight days before its cameras failed, probably from radiation damage to the vidicon targets.

None of this mattered. What mattered was the principle: you could see the weather from above. You could see the shape of a storm system in a single image, rather than inferring it from scattered ground observations. You could see cloud patterns over the ocean where there were no weather stations --- and the ocean covers seventy percent of the Earth. You could see the birth of a tropical cyclone days before it was detectable by any surface instrument. You could watch cold fronts march across continents. You could see fog banks along coastlines, snow cover over mountain ranges, jet stream cirrus at forty thousand feet. You could see the weather the way the weather actually exists --- as a three-dimensional fluid system wrapping the entire planet, not as a set of numbers from isolated observation points.

\begin{center}
$\bullet\quad\bullet\quad\bullet$
\end{center}

I went home that evening and sat on my back porch and looked up at the sky. It was partly cloudy --- the weather bureau's own forecast called for partly cloudy, and for once we were right. Cumulus clouds drifted across the sky at perhaps three thousand feet, their flat bases and rounded tops catching the late-afternoon sun. I had looked at clouds all my life. I had classified them, measured their bases with ceilometers and their tops with pilot reports, estimated their coverage in eighths of the sky --- the traditional oktas. I knew the difference between cumulus humilis and cumulus congestus, between cirrostratus and cirrocumulus, between a cumulonimbus and an altostratus opacus. I knew clouds.

But I had only ever seen them from below. I had seen their bottoms, their sides, their edges. I had never seen their tops. I had never seen them as they exist from above --- as a white surface against the dark Earth, a reflective layer between the ground and the sun, a texture that reveals the dynamics of the atmosphere the way ripples reveal the currents in a stream. TIROS-1 had given me the view from above, and the view from above was a different thing entirely.

I could hear music from a neighbor's window --- a piano, playing something I recognized as Rachmaninoff. The second piano concerto, I thought, though I was not sure. The music built slowly, the way weather builds --- a quiet theme in the lower register, gathering harmonics, adding voices, the dynamics swelling from \emph{piano} to \emph{forte} to \emph{fortissimo} in a great rolling crescendo that was not loud so much as \emph{complete}. A storm builds the same way. The initial disturbance is quiet: a surface low forms, the isobars contract slightly, the winds begin to organize. Then the convection starts, the moisture rises, the latent heat releases, the pressure drops further, the winds accelerate, and the system feeds on itself in a positive feedback loop that drives it from a modest disturbance to a raging typhoon in forty-eight hours. Rachmaninoff understood this. His music does not announce the storm --- it \emph{becomes} the storm, the listener swept up in the building energy until the resolution comes and everything is, for a moment, still.

The clouds above me drifted east. Tomorrow they would be somewhere else. A weather satellite overhead would see them in context --- part of a larger pattern, a brushstroke in a painting the size of a hemisphere. I could not see that pattern from my porch. But I had seen it in a photograph that morning, and I would never look at clouds the same way again.

\begin{center}
$\bullet\quad\bullet\quad\bullet$
\end{center}

\noindent\textit{The jellyfish pulses. The drum spins. The camera sees the Earth for one hundred degrees of each revolution, and in that brief window it captures what no human eye has ever seen from below: the shape of the weather, the spiral of the storm, the texture of the atmosphere laid bare from above. The first photograph was a revelation. The twenty-two thousand, nine hundred and fifty-second was a profession. Between the first and the last, the weather stopped being something we measured and became something we could see --- the eye of the storm, the eye of the satellite, the weather eye that changed how we understand the sky.}

\vspace{2em}
\begin{center}
\noindent\rule{0.3\textwidth}{0.4pt}
\end{center}
\vspace{1em}

\begin{center}
\textit{For Sergei Rachmaninoff (1873--1943), whose storms always resolve}\\[0.5em]
{\small NASA Mission Series v1.15 --- TIROS-1}\\
{\small Miles Tiberius Foxglove --- Mukilteo, WA}\\
{\small March 30, 2026}
\end{center}

\end{document}
